% 内角和定理证明：过 A 作 BC 平行线, 三角拼平角
\documentclass[tikz,border=4pt]{standalone}
\usepackage{ctex}
\usepackage{amsmath}
\usetikzlibrary{calc, arrows.meta}
\begin{document}
\begin{tikzpicture}[scale=1.0, thick]
  \coordinate (A) at (0, 2.4);
  \coordinate (B) at (-2.2, 0);
  \coordinate (C) at (2.6, 0);
  \draw (A) -- (B) -- (C) -- cycle;

  % 过 A 作 BC 平行线 (水平线，BC 是水平)
  \draw[red, dashed, -{Stealth[length=2.5mm]}] (A) -- (-2.5, 2.4);
  \draw[red, dashed, -{Stealth[length=2.5mm]}] (A) -- ( 2.5, 2.4);

  % 标三个角在 A 处
  \node at ($(A)+(-0.55, -0.15)$) {\footnotesize $\angle 1$};
  \node at ($(A)+( 0.55, -0.15)$) {\footnotesize $\angle 2$};
  \node at ($(A)+( 0,    -0.45)$) {\footnotesize $\angle A$};
  % 标 B, C 内角
  \node at ($(B)+( 0.4,   0.18)$) {\footnotesize $\angle B$};
  \node at ($(C)+(-0.5,   0.18)$) {\footnotesize $\angle C$};

  \node[above]       at (A) {$A$};
  \node[below left]  at (B) {$B$};
  \node[below right] at (C) {$C$};

  \node[red] at (-1.8, 2.65) {\footnotesize 过 $A$ 的 $BC$ 平行线};
\end{tikzpicture}
\end{document}
